% ---------- ¶13  related literature, three strands ----------
Our paper draws on three strands of work. It belongs to the literature that estimates
disclosure models structurally in order to assess the descriptiveness of classical disclosure
theory and to recover primitives such as information endowment, disclosure costs, and beliefs;
closest to us is \citet{bourveau_learning_2025}, who structurally examine whether the unraveling
result describes an unregulated nineteenth-century market and find that full disclosure is
reached only gradually. Where that literature and the experimental work on disclosure ask how far
receivers infer the worst from silence \citep{jin_is_2021, hagenbach_motivated_2025}, our extension
asks whether they attend to silence at all, drawing on evidence that limited attention changes how
disclosed and undisclosed information is priced \citep{hirshleifer_limited_2003, dellavigna_investor_2009}.
It is closely related to the extensive evidence on the capitalisation of energy labels in
housing, which documents consistent premia for efficient properties wherever labels have been
studied \citep{brounen_kok_2011, hyland_lyons_2013, fuerst_energy_2016, frondel_power_2018,
myers_mandatory_2022}; that literature establishes that labels are priced, whereas our question
is why their \emph{absence} is priced favourably. And it builds on work on information
asymmetry in real-estate markets, where the informational advantage of one side of the
transaction has been inferred from prices and time-on-market \citep{levitt_syverson_2008,
lizzeri_1999}.
